\documentclass[aspectratio=169]{beamer}

% --- PACKAGES ---
\usepackage{graphicx, bbm, amsmath, amsthm, amssymb}
\usepackage{booktabs}
\usepackage{array}
\usepackage{caption}
\usepackage{threeparttable}
\usepackage{subcaption}
\usepackage{xfp}
\usepackage{natbib}
\bibliographystyle{plainnat}

% --- THEME AND COLORS ---
\usetheme{moloch}
\setbeamercolor{background canvas}{bg=white}
\definecolor{cardinalred}{rgb}{0.549, 0.082, 0.082}
\setbeamercolor{title separator}{fg=cardinalred}
\setbeamercolor{progress bar in section page}{fg=cardinalred}
\setbeamercovered{invisible}
\setbeamertemplate{footline}[frame number]
\setbeamersize{text margin left=1.5em,text margin right=1em}

% --- SCALAR INPUTS ---
\input{../input/Tables/enoe_scalars}
\input{../input/Tables/homicide_scalars}
\input{../input/Tables/cocaine_scalars}
\newcommand{\FirstStageCoef}{ -0.0000008568 }
\newcommand{\FirstStageSe}{ 0.0000000689 }
\newcommand{\Fstat}{ 154.75 }
\newcommand{\RespikeCoef}{ 0.307 }
\newcommand{\HighCoef}{ 0.413 }
\newcommand{\MaleCoef}{ 0.495 }
\newcommand{\MunicipalHRSD}{ 0.87 }
\newcommand{\sdEffectRespike}{ 0.175 }
\newcommand{\sdEffectRespikeHigh}{ 0.235 }
\newcommand{\sdEffectRespikeMale}{ 0.281 }
\newcommand{\DistanceFirstStageEffect}{ 0.0020216688 }
\newcommand{\DistanceSd}{  192 }
\newcommand{\OLSCoef}{ -0.003 }
\newcommand{\OLSSe}{ 0.000 }
\newcommand{\TWFECoef}{ 0.044 }
\newcommand{\TWFESe}{ 0.084 }
\newcommand{\MainCoef}{ 0.067 }
\newcommand{\MainSe}{ 0.087 }
\newcommand{\WarCoef}{ -0.197 }
\newcommand{\WarSe}{ 0.565 }
\newcommand{\InterimCoef}{ -0.035 }
\newcommand{\InterimSe}{ 0.139 }
\newcommand{\RespikeSe}{ 0.137 }


% --- TITLE PAGE ---
\title{Crime and the Classroom}
\subtitle{How Local Violence Affects School Dropout Decisions in Mexico}
\author{Daniel Werner}
\institute{Stanford University}
\date{}

% ======================================================
\begin{document}

% --- TITLE SLIDE ---
\begin{frame}[plain]
\titlepage
\end{frame}

% ======================================================
\section{Introduction}

\begin{frame}{Mexico's violence epidemic}
\begin{columns}
\begin{column}{0.55\textwidth}
    \begin{itemize}
        \item Drug trafficking organizations (DTOs) generate \$19--29 billion USD annually
        \item Three distinct eras of violence since 2007:
        \begin{itemize}
            \item \textbf{War on Drugs} (2007--2012): Calder\'on's militarized offensive
            \item \textbf{Interim Period} (2013--2016): Strategic withdrawal, lower rates
            \item \textbf{Resurgence} (2017--2024): Historically high, persistent violence
        \end{itemize}
        \item 2016: Mexico was the world's \#2 deadliest conflict zone (behind Syria)
        \item In 2024: 35,600 persons reported missing
    \end{itemize}
\end{column}
\begin{column}{0.45\textwidth}
    \begin{figure}
        \centering
        \includegraphics[width=\textwidth]{../input/Figures/Homicide_TS.pdf}
    \end{figure}
\end{column}
\end{columns}
\end{frame}


\begin{frame}{Violence is geographically concentrated}
\begin{figure}
    \centering
    \includegraphics[width=0.78\textwidth]{../input/Figures/total_hom_map.pdf}
\end{figure}
\vspace{-0.3em}
{\footnotesize Average municipal homicide rates 2007--2024. Darkest municipalities are along the Pacific coast.}
\end{frame}


\begin{frame}{Research question}

    \begin{block}{}
        \centering
        Does an increase in the local homicide rate cause young people to drop out of school?
    \end{block}

\vspace{0.5cm}
\pause

The effect is \textbf{theoretically ambiguous}:

\vspace{0.3cm}
\begin{columns}
\begin{column}{0.47\textwidth}
    \textbf{Violence $\uparrow$ dropout:}
    \begin{itemize}
        \item Fear reduces perceived safety
        \item Children supplement lost household income
        \item DTOs recruit youth as enforcers, runners, lookouts
    \end{itemize}
\end{column}
\begin{column}{0.47\textwidth}
    \textbf{Violence $\downarrow$ dropout:}
    \begin{itemize}
        \item Schools function as safe havens
        \item Higher enforcement risk deters participation
        \item Social stigma of drug trade discourages recruitment
    \end{itemize}
\end{column}
\end{columns}

\vspace{0.5cm}
\pause
$\implies$ This is an \textbf{empirical question}

\end{frame}


\begin{frame}{This paper}

\textbf{Research question}: What is the causal effect of drug-trafficking violence on school dropout in Mexico?

\pause
\vspace{0.5cm}

\textbf{Approach}:
\begin{itemize}
    \item Individual-level panel data from ENOE (2007--2024): \totalKids\ children
    \item Construct high-frequency monthly dropout indicator
    \item Instrumental variables: Colombian cocaine seizures $\times$ distance to Pacific coast
    \item Exploit three eras of violence and enforcement policy
\end{itemize}

\pause
\vspace{0.5cm}

\textbf{Preview of findings}:
\begin{itemize}
    \item Pooled IV estimate: small positive effect, imprecise
    \item \textbf{Resurgence period (2017--2024)}: large, positive, statistically significant
    \item Effects concentrated among \textbf{high school--aged males}
    \item Evidence consistent with \textbf{DTO recruitment mechanism}
\end{itemize}

\end{frame}


\begin{frame}{Contributions}
\begin{enumerate}
    \item \textbf{Identification}: IV strategy (shift-share) addresses reverse causality and omitted variable bias present in prior TWFE estimates

    \vspace{0.4cm}

    \item \textbf{Sample period}: Covers 2007--2024 --- the longest panel in this literature; most prior work covers only 2007--2012

    \vspace{0.4cm}

    \item \textbf{Outcome}: First monthly, individual-level dropout measure for Mexico; avoids institutional reporting incentives in administrative data

    \vspace{0.4cm}

    \item \textbf{New finding}: Evidence of a structural break under AMLO's permissive enforcement regime, consistent with heightened DTO recruitment
\end{enumerate}
\end{frame}


% ======================================================
\section{Data}

\begin{frame}{Data}
\begin{enumerate}
    \item ENOE: National Survey of Occupation and Employment
    \item Homicides: INEGI official death records
    \item Colombian cocaine seizures: Colombian Ministry of Defense
    \item Geographic data: INEGI shapefiles
\end{enumerate}
\end{frame}


\begin{frame}{ENOE: rotating individual panel}
\begin{columns}
\begin{column}{0.55\textwidth}
    \begin{itemize}
        \item Nationally representative rotating panel, 5 consecutive quarters per individual
        \item Covers 2007--2024; one quarter missing (COVID, Q2 2020)
        \item Keep only individuals completing all 5 surveys
        \item \textbf{Sample}: \totalKids\ children aged 12--18 in urban municipalities ($\geq$15,000 pop.)
        \item \textbf{Dropout}: transition from enrolled $\rightarrow$ not enrolled between two survey waves
        \begin{itemize}
            \item Sidesteps institutional reporting incentives
        \end{itemize}
    \end{itemize}
\end{column}
\begin{column}{0.45\textwidth}
    \begin{itemize}
        \item Dropout rates by school level:
        \begin{itemize}
            \item Secondary: \dRateSecondary\%
            \item High school: \dRateHigh\%
        \end{itemize}
        \item Discontinuity at age 15 = minimum legal working age
        \item Primary-age children (6--11) excluded:
        \begin{itemize}
            \item Negligible dropout
            \item Theoretical channels don't apply
        \end{itemize}
    \end{itemize}
\end{column}
\end{columns}
\end{frame}


\begin{frame}{Enrollment rates have been roughly flat since 2015}
\begin{figure}
    \centering
    \includegraphics[width=0.82\textwidth]{../input/Figures/enrollment_TS.pdf}
\end{figure}
\vspace{-0.3em}
{\footnotesize Average monthly enrollment rate among secondary and high school students (ENOE).}
\end{frame}


\begin{frame}{Homicide data: three eras of violence}
\begin{figure}
    \centering
    \includegraphics[width=0.82\textwidth]{../input/Figures/Homicide_TS.pdf}
\end{figure}
\vspace{-0.5em}
{\footnotesize National and average municipal homicide rates per 10,000 inhabitants (INEGI). Vertical lines mark period boundaries.}
\end{frame}


\begin{frame}{Colombian cocaine seizures: the supply shock}
\begin{columns}
\begin{column}{0.55\textwidth}
    \begin{itemize}
        \item Large coastal seizures only ($\geq$500 kg)
        \item \nSeizeWholeSamp\ seizure events/month on average
        \item Average of \kgSeizeWholeSamp\ metric tonnes/month
        \item Standard deviation: \sdSeizeWholeSamp\ metric tonnes
        \item High correlation with Mexican homicide rates
        \item Variation is mostly idiosyncratic (Colombian enforcement operations)
    \end{itemize}
\end{column}
\begin{column}{0.45\textwidth}
    \begin{figure}
        \centering
        \includegraphics[width=\textwidth]{../input/Figures/Seizure_TS.pdf}
    \end{figure}
\end{column}
\end{columns}
\end{frame}


% ======================================================
\section{Empirical Strategy}

\begin{frame}{Empirical strategy}

\textbf{Challenge}: Violence is endogenous
\begin{itemize}
    \item Reverse causality: dropouts may themselves affect violence
    \item Omitted variable bias: unobserved factors drive both
\end{itemize}

\vspace{0.5cm}
\pause

\textbf{Solution}: Instrumental Variables (IV)
\begin{itemize}
    \item Instrument: $\underbrace{d_i}_{\text{distance to Pacific}} \times \underbrace{c_{j,t}}_{\text{Colombian seizures}}$
    \item Shift-share design, following \citet{dubevargas}
\end{itemize}

\vspace{0.5cm}
\pause

\textbf{Identifying assumption}: Colombian seizures $\times$ Pacific distance affect dropout \emph{only} through changes in local homicide rates

\end{frame}


\begin{frame}{The instrument: two components}

\textbf{1. Colombian cocaine seizures} ($c_{j,t}$): supply shocks
\begin{itemize}
    \item Large Colombian seizures reduce cocaine supply
    \item Short-run demand is inelastic $\Rightarrow$ retail prices rise
    \item Mexican DTOs compete violently for control of trafficking routes
    \item Driven by Colombian enforcement operations: orthogonal to Mexico
\end{itemize}

\vspace{0.5cm}
\pause

\textbf{2. Distance to Pacific coast} ($d_i$): exposure
\begin{itemize}
    \item $\approx$70\% of Colombian cocaine enters Mexico via Pacific routes
    \item Closer to Pacific $=$ more important within the trafficking network
    \item Time-invariant: identification from differential exposure to common shocks
    \item Consistent with Figure: highest homicide rates along Pacific corridor
\end{itemize}

\end{frame}


\begin{frame}{Model}

\textbf{First stage}:
\small
\begin{equation*}
    \ln(1 + h_{i,j,t}) = \gamma_1 + \gamma_2 \bigl[ d_i \times c_{j,t} \bigr] + \mathbf{X}_{i,j,t}'\Gamma + \mathbf{W}_{s,i,j,t}'\Theta + \delta_{j,t} + m_s + \nu_{i,j,t}
\end{equation*}

\normalsize
\vspace{0.4cm}
\pause

\textbf{Second stage}:
\small
\begin{equation*}
    y_{s,i,j,t} = \beta_1 + \beta_2\,\widehat{\ln(1 + h_{i,j,t})} + \mathbf{X}_{i,j,t}'\Gamma + \mathbf{W}_{s,i,j,t}'\Theta + \delta_{j,t} + m_s + \varepsilon_{s,i,j,t}
\end{equation*}

\normalsize
\vspace{0.4cm}

where:
\begin{itemize}
    \item $y_{s,i,j,t}$ = 1 if student $s$ drops out in municipality $i$, year $j$, month $t$
    \item $\mathbf{X}_{i,j,t}$ = municipal controls (employment, income, population, lags of violence)
    \item $\mathbf{W}_{s,i,j,t}$ = individual/household controls
    \item $m_s$ = individual fixed effects; $\delta_{j,t}$ = month--by--year fixed effects
    \item Standard errors clustered at the individual level
\end{itemize}

\end{frame}


\begin{frame}{Instrument validity: uncorrelated with Mexican enforcement}

\begin{center}
\scalebox{0.72}{%
\begin{threeparttable}
\begin{tabular}{l c c c c c}
\hline\hline
 & Marijuana & Heroin & Meth & Opium & Total \\
 & (1) & (2) & (3) & (4) & (5) \\
\hline
Colombian cocaine seizures (metric tonnes)
 & $-4.220^{}$ & $ 0.001^{*}$ & $-0.034^{}$ & $-0.004^{}$ & $-4.262^{}$  \\
 & ( 3.448) & ( 0.001) & ( 0.053) & ( 0.003) & ( 3.467)  \\
\hline
\hline
Observations
 & 12 & 12 & 12 & 12 & 12 \\
\hline\hline
\end{tabular}

\end{threeparttable}
}
\end{center}
\vspace{0.3em}
{\footnotesize Colombian seizures have \textbf{no significant correlation} with Mexican drug interdiction efforts across any substance (2007--2018). Annual data; 12 observations.}

\end{frame}


\begin{frame}{Instrument validity: uncorrelated with economic outcomes}

\begin{center}
\scalebox{0.72}{%
\begin{threeparttable}
\input{../input/Tables/placebo_seizures}
\end{threeparttable}
}
\end{center}
\vspace{0.3em}
{\footnotesize Colombian seizures are uncorrelated with municipal employment rates, average income, weekly hours, household income, or national GDP.}

\end{frame}


% ======================================================
\section{Results}

\begin{frame}{Main results}

\begin{center}
\scalebox{0.62}{%
\begin{threeparttable}
\begin{tabular}{l c c c c c c c }
\hline\hline
 {\small \textit{Outcome: Dropout}} & \shortstack{OLS \\ (1)} & \shortstack{TWFE  \\ (2)} & \shortstack{Minimal \\ controls \\ (3)} & \shortstack{Primary \\ specification \\ (4)} & \shortstack{War on drugs \\ 2007-2012 \\ (5)} & \shortstack{Interim \\ 2013-2016 \\ (6)} & \shortstack{Resurgence \\ 2017-2024 \\  (7)} \\
\hline
Homicides per 10,000
 & $-0.003^{***}$ & $ 0.001^{}$ & $ 0.044^{}$ & $ 0.067^{}$ & $-0.197^{}$ & $-0.035^{}$ & $ 0.307^{**}$  \\
 & ( 0.000) & ( 0.001) & ( 0.084) & ( 0.087) & ( 0.565) & ( 0.139) & ( 0.137)  \\
\hline
Mean of school enrollment
 & $ 0.842$ & $ 0.843$ & $ 0.843$ & $ 0.843$ & $ 0.821$ & $ 0.857$ & $ 0.860$ \\
Kleibergen-Paap F-stat
 & –- & –- & 164.06 & 154.75 &   3.96 &  64.02 &  65.44 \\
\hline
Observations
 & 2,032,058 & 1,857,143 & 1,973,879 & 1,857,143 &   742,496 &   507,571 &   578,125 \\
\hline\hline
\end{tabular}

\end{threeparttable}
}
\end{center}

\end{frame}


\begin{frame}{Three key findings from baseline estimates}

\begin{enumerate}

    \item \textbf{OLS and TWFE} yield small, negative estimates --- likely reflecting selection bias and time-invariant confounders

    \vspace{0.4cm}

    \item \textbf{Pooled IV}: positive point estimates --- reversal from OLS suggests endogeneity was masking a positive effect; but pooled estimate is imprecise

    \vspace{0.4cm}

    \item \textbf{Time heterogeneity} (the central finding):
    \begin{itemize}
        \item War on Drugs (2007--12): negative, imprecise; weak first stage
        \item Interim Period (2013--16): negative, tight around zero; strong first stage
        \item \textbf{Resurgence (2017--24): large, positive, significant at 5\%}
    \end{itemize}

\end{enumerate}

\vspace{0.3cm}
\pause
Kleibergen--Paap F-stat (pooled): \Fstat\ $\gg$ conventional thresholds

\end{frame}


\begin{frame}{Resurgence period: effects concentrated in older males}

\begin{center}
\scalebox{0.72}{%
\begin{threeparttable}
\input{../input/Tables/subgroup_iv_respike}
\end{threeparttable}
}
\end{center}

\end{frame}


\begin{frame}{Magnitude of the Resurgence effect}

A one--standard deviation increase in the municipal homicide rate (\MunicipalHRSD\ per 10,000):

\vspace{0.4cm}

\begin{itemize}
    \item \textbf{All students (12--18)}: increases dropout probability by $\approx$ {\fpeval{round(\sdEffectRespike * 100,2)}} pp
    \begin{itemize}
        \item Mean dropout rate $\approx$ 7.4\%
    \end{itemize}

    \vspace{0.3cm}

    \item \textbf{High school students (15--18)}: $\approx$ {\fpeval{round(\sdEffectRespikeHigh * 100,2)}} pp
    \begin{itemize}
        \item Mean $\approx$ 9.5\%
    \end{itemize}

    \vspace{0.3cm}

    \item \textbf{High school males}: $\approx$ {\fpeval{round(\sdEffectRespikeMale * 100,2)}} pp
    \begin{itemize}
        \item Mean $\approx$ 8\%
    \end{itemize}
\end{itemize}

\vspace{0.4cm}
\pause

Results are \textbf{robust} to quarterly aggregation (Appendix Table)

\end{frame}


% ======================================================
\section{Mechanism}

\begin{frame}{Mechanism: DTO recruitment, not safety fears}

\textbf{If the mechanism were fear of physical danger}, we would expect:
\begin{itemize}
    \item Effects concentrated among \textit{younger} and \textit{female} students (more vulnerable)
\end{itemize}

\vspace{0.4cm}
\pause

\textbf{But the data show the opposite}:
\begin{itemize}
    \item Larger effects for \textit{older} students (high school $>$ secondary)
    \item Concentrated in \textit{males}; null effects for females
\end{itemize}

\vspace{0.4cm}
\pause

\textbf{Consistent with DTO recruitment}:
\begin{itemize}
    \item DTOs recruit older adolescent males as enforcers, distributors, operatives
    \item Younger teens serve in peripheral roles (lookouts) compatible with school
    \item Older teens accumulate ``criminal capital'' $\rightarrow$ full-time criminal careers \citep{Sviatschi2022}
    \item DTO recruitment overwhelmingly targets young men
\end{itemize}

\end{frame}


\begin{frame}{Why the Resurgence period specifically?}

Effect is large in 2017--2024 but not before. Why?

\vspace{0.4cm}
\pause

\textbf{AMLO's ``abrazos, no balazos'' (hugs, not bullets)}:
\begin{itemize}
    \item Explicitly non-confrontational stance toward organized crime
    \item Reduced perceived risk of imprisonment or death from cartel membership
    \item Implicitly reduced the social stigma of participating in the drug trade
    \item Under Calder\'on: intense crackdowns $\Rightarrow$ high mortality risk of cartel membership
    \item Under AMLO: weak deterrence $\Rightarrow$ opportunity cost of school falls
\end{itemize}

\vspace{0.4cm}
\pause

\textbf{Mexico's National Registry of Missing Persons}:
\begin{itemize}
    \item 3,000 missing in 2007 $\rightarrow$ 11,000 in 2012 $\rightarrow$ 35,600 in 2024
    \item Age-gender distribution: males peak at 15--19 and remain elevated through middle age
    \item Precisely the age range where dropout effects are largest
\end{itemize}

\end{frame}


% ======================================================
\section{Robustness}

\begin{frame}{Robustness checks}

\begin{center}
\scalebox{0.62}{%
\begin{threeparttable}
\input{../input/Tables/robustness_iv}
\end{threeparttable}
}
\end{center}

\end{frame}


\begin{frame}{Robustness: summary}

Positive Resurgence estimate is robust to:

\begin{itemize}
    \item Including survey attriters (removing 5-quarter balance requirement)
    \item Lowering population threshold to 10,000
    \item Using all Colombian cocaine seizures (not just large coastal events)
    \item Dropping 2020 (COVID)
    \item Using school enrollment as outcome (sign reverses, as expected)
    \item Quarterly aggregation (separate Appendix)
\end{itemize}

\vspace{0.4cm}
\pause

Violence does \textbf{not} predict survey attrition:
\begin{itemize}
    \item Small, statistically insignificant coefficients across all periods
    \item Only exception: slightly negative during the Interim period
\end{itemize}

\end{frame}


% ======================================================
\section{Conclusion}

\begin{frame}{Conclusion}

\textbf{Question}: Does drug-trafficking violence cause school dropout in Mexico?

\vspace{0.5cm}

\textbf{Method}: IV with Colombian cocaine seizures $\times$ distance to Pacific coast; individual panel 2007--2024

\vspace{0.5cm}
\pause

\textbf{Findings}:
\begin{itemize}
    \item Pooled effect is small and positive (larger than TWFE estimates in the literature)
    \item \textbf{Resurgence period (2017--2024)}: large, positive, and significant at 5\%
    \item Effects are concentrated among \textbf{older males}
\end{itemize}

\vspace{0.3cm}
\pause

\textbf{Interpretation}: Consistent with DTO recruitment when enforcement is permissive. Not consistent with fear-based or income-substitution mechanisms.

\vspace{0.3cm}
\pause

\textbf{Policy implication}: Scholarship programs targeting high school males in high-violence areas may reduce recruitment incentives of criminal organizations

\end{frame}


% ======================================================
% References (hidden from presentation; needed for \citet{} to resolve)
\begin{frame}[allowframebreaks]{References}
\bibliography{../../4_paper/source/refs}
\end{frame}

\end{document}
